\section{Cluster truncation on a chain of interaction labels}
\label{sec:interval_cluster_truncation}

\subsection{Required labels with an empty forbidden set}

For a finite set $U$ of required interaction labels, let
\begin{align}
    R_U^H(\beta)
     & =\sum_{D\subseteq U}(-1)^{|D|}
    e^{-\beta(H-H_D)},
     & H_D                            & =\sum_{i\in D}H_i.
    \label{eq:mandatory_cluster_part}
\end{align}
This is the cluster part~(\ref{eq:hastings_cluster_part}) with the forbidden
set equal to $\varnothing$. The normalization below is an auxiliary
estimate for the interaction-label expansion related to the argument
of~\cite[lines 1270--1291 and 1935--1968]{Arad2016RigorousRG}.
It is distinct from the printed spatial approximation estimate, whose
obstructions are recorded in~\cite{gap:alvv16_cluster_approximation}.

\subsection{A normalized tail for an interaction-label cutoff}

The cutoff in this subsection counts interaction labels. It differs from
both physical-site support and the unrestricted estimate discussed
in~\cite{gap:alvv16_cluster_approximation}. No positivity of the individual
interactions is imposed.

\subsection{Physical factors at a cut of the label path}

Let $K$ be a finite interaction-label set with graph $G$, let $H_i$ be
complex square matrices, and let $E_A$ be their exact-label operators.
For a fixed criterion of large components, write
\begin{align}
    F(B)=\sum_{\substack{A\subseteq B \\
            A\text{ has no large component}}}E_A.
    \label{eq:retained_allowed_labels}
\end{align}
Every operator in (\ref{eq:retained_allowed_labels}) is a finite sum of
the existing exact-label coefficients.

\begin{lemma}[Physical nearest-neighbour crossing decomposition]
    \label{thm:physical_cluster_crossing_decomposition}
    \lean{TNLean.Lattice.ClusterExpansion.overlapGraph_nearestNeighbor_eq_pathGraph,
        TNLean.Lattice.ClusterExpansion.commute_nearestNeighbor_of_not_adj,
        ClusterExpansion.truncatedExp_pathGraph_eq_crossing_components,
        TNLean.Lattice.ClusterExpansion.truncatedExp_nearestNeighbor_eq_crossing_components}
    \leanok
    
    Let $H_i$, $0\leq i<N$, be operators on $N+1$ sites, with
    $H_i$ supported on the designated pair $\{i,i+1\}$. Fix an integer
    $r\geq1$ and a label $p$, corresponding to the physical cut after
    site $p$. In (\ref{eq:retained_allowed_labels}), retain supports
    whose label components have fewer than $r$ labels. For every real
    $\beta$, their actual selected expansion satisfies
    \begin{align}
        M_r=F(L_p)F(R_p)+
        \sum_{\substack{C\text{ an interval of labels} \\p\in C,\ |C|<r}}
        E_C\bigl(F(L(C))F(R(C))\bigr).
        \label{eq:physical_cluster_crossing_decomposition}
    \end{align}
    All factors are the finite sums in
    (\ref{eq:retained_allowed_labels}). Neither Hermiticity nor
    positivity of the interactions is needed for this identity.
    When $r=1$ the sum over crossing components is empty.
\end{lemma}
\begin{proof}\leanok
    
    The overlap graph of the designated pairs $\{i,i+1\}$ is exactly
    the path on the interaction labels. Nonadjacent terms act on
    disjoint physical sites and commute. Split the retained supports
    according to whether they contain $p$. The first family is a
    product of the two sums on $L_p$ and $R_p$. In the second family,
    use the unique component containing $p$ and split its exterior
    into $L(C)$ and $R(C)$.
\end{proof}

\begin{lemma}[Physical support of retained side sums]
    \label{thm:retained_cluster_physical_support}
    \lean{TNLean.Lattice.ClusterExpansion.isSupportedOn_sum_retained_clusterOp,
        TNLean.Lattice.ClusterExpansion.isSupportedOn_retained_nearestNeighbor_interval,
        TNLean.Lattice.ClusterExpansion.isSupportedOn_retained_nearestNeighbor_left,
        TNLean.Lattice.ClusterExpansion.isSupportedOn_retained_nearestNeighbor_right}
    \leanok
    \uses{thm:hastings_cluster_operator_support}
    If each $H_i$ acts on a designated set $X_i$, then $F(B)$ acts
    only on $\bigcup_{i\in B}X_i$. In the nearest-neighbour case,
    suppose $a\leq i$ and $i+1<b$ for every $i\in B$. Then $F(B)$
    acts only on the physical interval $[a,b)$.
    In particular, labels $i<p$ give a factor acting only on sites
    at most $p$, and labels $p<i$ give a factor acting only on sites
    strictly greater than $p$. These conclusions apply to all the
    side factors in (\ref{eq:physical_cluster_crossing_decomposition}).
\end{lemma}
\begin{proof}\leanok
    \uses{thm:hastings_cluster_operator_support}
    Each summand $E_A$, $A\subseteq B$, is supported on
    $\bigcup_{i\in A}X_i\subseteq\bigcup_{i\in B}X_i$.
    The support algebra is closed under finite sums. For a bond
    label $i$, both endpoints $i$ and $i+1$ lie in the stated physical
    interval. The two one-sided conclusions follow from the strict
    inequalities relative to $p$.
\end{proof}

These identities use the corrected interaction-label cutoff. A retained
component of $s$ labels is designated on $s+1$ physical sites; the cutoff
is not the source's selector based on the minimal physical support of a
word product. The normalization and error estimate of
the normalized interaction-label tail on a path
\cite[``Normalized interaction-label tail on a path'']{QICLean2026Blueprint} are unchanged. The
Schmidt-rank count and a matrix product operator representation are
established next.

\begin{lemma}[Two-sided rank of a crossing interval]
    \label{thm:physical_cluster_interval_rank}
    \lean{TNLean.Lattice.ClusterExpansion.card_biUnion_nearestNeighbor_le,
        TNLean.Lattice.ClusterExpansion.operatorSchmidtRank_clusterOp_nearestNeighbor_le}
    \leanok
    \uses{thm:hastings_cluster_operator_support}
    Let every site have dimension $q\geq1$. If $C$ is a connected
    interval of nearest-neighbour labels, its designated physical
    support has at most $|C|+1$ sites. If $|C|<r$, then at every
    physical cut the exact-label coefficient $E_C$ satisfies
    \begin{align}
        \operatorname{OSR}(E_C)\leq q^r.\notag
    \end{align}
\end{lemma}
\begin{proof}\leanok
    \uses{thm:hastings_cluster_operator_support}
    Write $C=[a,b]$. All its bonds act within the physical interval
    $[a,b+1]$, which has $|C|+1\leq r$ sites. On a designated support
    $X$, the matrix-unit bound on the two sides of a cut $A$ gives
    \begin{align}
        \operatorname{OSR}_A(T)
        \leq\min\{q^{2|X\cap A|},q^{2|X\setminus A|}\}
        \leq q^{|X|}.\notag
    \end{align}
    The two parts partition $X$, so twice the smaller cardinality
    is at most $|X|$. Apply this to the proved support of $E_C$.
\end{proof}

\begin{lemma}[All cut ranks of the label truncation]
    \label{thm:physical_cluster_all_cut_rank}
    \lean{TNLean.Lattice.ClusterExpansion.operatorSchmidtRank_truncatedExp_nearestNeighbor_le_count,
        TNLean.Lattice.ClusterExpansion.operatorSchmidtRank_truncatedExp_nearestNeighbor_le,
        TNLean.MPS.AreaLaw.operatorSchmidtRank_cluster_truncation_all_cuts_le}
    \leanok
    \uses{thm:physical_cluster_crossing_decomposition, thm:retained_cluster_physical_support, thm:physical_cluster_interval_rank}
    Let $r\geq1$, $q\geq1$, and let the interactions have the actual
    designated nearest-neighbour supports. For every real $\beta$,
    at an interior cut the selected operator satisfies
    \begin{align}
        \operatorname{OSR}(M_r)
        \leq1+(r-1)^2q^r\leq r^2q^r.\notag
    \end{align}
    The bound $r^2q^r$ holds at every chain cut, including the two
    endpoint cuts. No Hermiticity or positivity is required.
\end{lemma}
\begin{proof}\leanok
    \uses{thm:physical_cluster_crossing_decomposition, thm:retained_cluster_physical_support, thm:physical_cluster_interval_rank}
    In the exact crossing decomposition, the product with no
    crossing component has rank at most one. Each crossing term
    has rank at most $q^r$: its central factor has this rank, and
    both exterior factors act on only one side. Subadditivity and
    submultiplicativity of operator Schmidt rank give the first
    inequality. For $r\geq1$, $1+(r-1)^2\leq r^2$ and $q^r\geq1$.
    An endpoint cut has a one-dimensional configuration space on
    one side. If $q=1$ the whole operator space is one-dimensional;
    if $r=1$ the crossing family is empty. These observations also
    cover a chain with one site and no interactions.
\end{proof}

\begin{lemma}[Operator tensors from the actual cut ranks]
    \label{thm:operator_chain_from_all_cut_ranks}
    \lean{MPOTensor.rank_pairedOperator_cutMatrix_eq,
        MPOTensor.exists_traceChainOperator_of_operatorSchmidtRank_le}
    \leanok
    Let $T$ be an operator on a chain of $L\geq1$ sites of local
    dimension $q\geq0$. Let $D\geq1$ and suppose every chain cut
    has operator Schmidt rank at most $D$. Then there exist
    site-dependent matrices $A_v^{s,t}\in\mathbb C^{D\times D}$ such that
    \begin{align}
        T_{\sigma,\tau}
        =\operatorname{tr}\bigl(A_0^{\sigma_0,\tau_0}\cdots
        A_{L-1}^{\sigma_{L-1},\tau_{L-1}}\bigr).\notag
    \end{align}
\end{lemma}
\begin{proof}\leanok
    Pair the two physical indices at each site. At every cut, the
    resulting coefficient matrix is a row and column permutation
    of the realigned operator matrix, so the two ranks agree.
    The open-chain MPS construction therefore supplies actual
    rectangular tensors of internal dimension at most $D$.
    Pad these tensors to square matrices, retaining the coefficient
    because $L\geq1$. Their scalar closing factor is multiplied
    into the first tensor. This constructs the stated tensors.
\end{proof}

\begin{theorem}[Actual matrix product operator for the corrected truncation]
    \label{thm:physical_cluster_label_mpo}
    \lean{TNLean.MPS.AreaLaw.exists_mpo_cluster_truncation,
        TNLean.MPS.AreaLaw.exists_mpo_cluster_approximation}
    \leanok
    \uses{thm:physical_cluster_all_cut_rank, thm:operator_chain_from_all_cut_ranks}
    Let $r,q\geq1$, and let $H=\sum_{i=0}^{N-1}H_i$ on $N+1$ sites,
    with each $H_i$ supported on $\{i,i+1\}$. For every real $\beta$,
    the actual interaction-label truncation $M_r$ has a site-dependent
    matrix product operator representation of bond dimension $r^2q^r$.
    If, in addition, the terms are Hermitian, $\|H_i\|\leq1$,
    $H\geq\lambda I$, $\beta\geq0$, and $b=2(e^\beta-1)\leq1$, the
    same represented operator satisfies
    \begin{align}
        e^{\beta\lambda}\|e^{-\beta H}-M_r\|
        \leq e^{N^2b^r}-1.\notag
    \end{align}
\end{theorem}
\begin{proof}\leanok
    \uses{thm:physical_cluster_all_cut_rank, thm:operator_chain_from_all_cut_ranks}
    The chain has $N+1\geq1$ sites and $r^2q^r\geq1$. Apply the
    constructed all-cut rank estimate and then the preceding tensor
    construction. Substitute the resulting exact representation in
    the already proved normalized error bound.
\end{proof}

This is a corrected interaction-label theorem related to
\cite[Theorems on cluster approximation and representation, lines 1270--1303]{Arad2016RigorousRG}.
The printed physical-support approximation remains false; its
counterexamples and the distinction between label count and minimal
physical support remain in \cite{gap:alvv16_cluster_approximation}.
The result asserts the existence of the tensors of the actual finite
selected operator. It does not assert the source's computational
running-time bound.
